\section*{Chapter \ref{ch:mx:the-empirics-of-price-impact} --- The Empirics of Price Impact}

\begin{solution}{exo:mx:the-empirics-of-price-impact:1}
$0.8\times0.02\times\sqrt{0.01}=0.0016$, 16 basis points; $0.8\times0.02\times\sqrt{0.04}=0.0032$, 32 basis points.
\end{solution}

\begin{solution}{exo:mx:the-empirics-of-price-impact:2}
About $0.6\times2.1\approx1.3$ ticks. Small imbalances meet the thin levels near the price and move it proportionally; large ones reach deeper levels and
wake the fundamentalists, so the relation bends.
\end{solution}

\begin{solution}{exo:mx:the-empirics-of-price-impact:3}
All of it, within its error: $-1.5$ ticks with a standard error of 1.5. The fundamentalists trade toward the hidden value, which the metaorder did not
move, and pull the price back.
\end{solution}

\begin{solution}{exo:mx:the-empirics-of-price-impact:4}
The forecast drift is $a\sigma\sqrt T$ and $T$ is proportional to $Q/V$ at a fixed participation rate, so the drift is proportional to $\sigma\sqrt{Q/V}$,
the same shape as the impact: it adds to the prefactor ($0.8+a$ times the share timed, roughly) and leaves the slope.
\end{solution}

\begin{solution}{exo:mx:the-empirics-of-price-impact:5}
Orders on one stock share its volatility, its liquidity and its days, so their errors are correlated; resampling single orders treats them as independent
and gives intervals too narrow.
\end{solution}

\begin{solution}{exo:mx:the-empirics-of-price-impact:6}
Its orders were recorded in a market without ours: its providers keep posting where they posted and its takers keep taking what they took, whatever we
do. A backtest on replayed data has only the mechanical cost of the book we consumed, no response to it, and so underestimates the cost of large orders.
\end{solution}

\begin{solution}{exo:mx:the-empirics-of-price-impact:7}
Without timed orders the exponent falls to 0.28 (prefactor with the exponent fixed 0.81); with them, 0.20 (1.24). Over a tenth of a day the noise
$\sigma\sqrt T$ is the same for every order and swamps the small ones, whose bin means are dominated by noise; the timed orders' forecast, now the same
for every size, flattens the curve further.
\end{solution}

\begin{solution}{exo:mx:the-empirics-of-price-impact:8}
The exponent can be right while the level is not: in the simulation, 30\% of orders timed by a forecast left the exponent at 0.50 and raised the prefactor
by 62\%. Check the prefactor against orders without a forecast, or remove the forecast order by order.
\end{solution}

\begin{solution}{pb:mx:the-empirics-of-price-impact:1}
\textbf{1.} Providers placing orders relative to the current book, pre-drawn noise orders, fundamentalists trading toward a hidden value; the providers and
fundamentalists react.
\textbf{2.} 0.27, 0.58, 1.12 and 2.34 ticks for 1, 2--3, 4--7 and 8+ lots; after fifty orders 0.81, 0.23, 0.68 and 0.46.
\textbf{3.} The informed fundamentalists send one-lot orders, and the price moves toward the value after them.
\textbf{4.} 2.1 ticks.
\textbf{5.} The session with the metaorder less the same session without it, with the same noise orders and value jumps: the exogenous noise cancels.
\textbf{6.} No difference at any horizon: the replayed providers never learn of the order.
\textbf{7.} 6.3, 7.8 and $-1.5$ ticks.
\textbf{8.} The price differences after the first child grow with the providers' different choices; standard errors of 1 to 3 ticks against impacts of a
few ticks.
\textbf{9.} $I=Y\sigma(Q/V)^\delta$, $\delta\approx\tfrac12$; 16 and 32 basis points.
\textbf{10.} Mean impact over volatility within size bins, a log-log fit on the bin means, a bootstrap over stocks; single impacts can be negative.
\textbf{11.} $\delta=0.48$ (0.39 to 0.56), $Y=0.85$ (0.55 to 1.66); 0.81 with the exponent fixed.
\textbf{12.} A V-shaped average book vanishing at the price, and a square root fitting two orders of magnitude of size with a logarithm fitting five.
\textbf{13.} The price move the order's owner expected, counted as impact.
\textbf{14.} $\delta=0.50$ (0.44 to 0.55), $Y=1.42$ (1.01 to 2.04), 1.31 with the exponent fixed.
\textbf{15.} The forecast's horizon is the duration, which grows with size, so the drift has the impact's shape.
\textbf{16.} Subtract each timed order's forecast; the regression of impact on the forecast gives 1.54, because larger orders carry both.
\textbf{17.} The exponent falls to 0.28 without timed orders and 0.20 with them: noise and the size-independent forecast flatten the curve.
\textbf{18.} Selection of executed orders, overlapping metaorders, replayed backgrounds (and noise).
\textbf{19.} \emph{Named result}: on 2,000 simulated metaorders with a true law of 0.8 and one half, the fit gives $\delta=0.48$ (0.39 to 0.56) and $Y=0.85$ (0.55
to 1.66); with 30\% of orders timed by a forecast, $\delta=0.50$ and $Y=1.42$, the prefactor 62\% too high with the exponent unchanged.
\textbf{20.} A square root with the prefactor fitted on orders without a forecast, or with each order's forecast removed.
\end{solution}

\begin{solution}{iq:mx:the-empirics-of-price-impact:1}
$I=Y\sigma\sqrt{Q/V}$: the price move of a metaorder as a fraction of the price, $Y$ of order one, $\sigma$ daily volatility, $Q$ the order's size and $V$
daily volume; it holds on average, for orders that are a fraction of daily volume, whatever the execution schedule within broad limits.
\iqlookfor{Each term; what is averaged; the range of validity.}
\end{solution}

\begin{solution}{iq:mx:the-empirics-of-price-impact:2}
From start to last fill against the arrival mid, normalised by volatility, binned by size over volume, fitted in logs, clustered by stock and day. Biases:
forecasts in timing, selection of executed orders, overlap with other orders, noise, and the horizon chosen.
\iqlookfor{Normalisation; binning; clustering; \omterm{def:mx:the-empirics-of-price-impact:alpha}{alpha contamination}.}
\end{solution}

\begin{solution}{iq:mx:the-empirics-of-price-impact:3}
Liquidity near the price is small and grows away from it (a V-shaped latent book): consuming $Q$ requires a move whose square grows with $Q$. Also,
large orders are split and the book refills between children.
\iqlookfor{Latent liquidity; splitting; refill.}
\end{solution}

\begin{solution}{iq:mx:the-empirics-of-price-impact:4}
The replayed market cannot respond: only the book the order consumed is charged, and nobody moves prices against it. Use an impact model or a reactive
simulator.
\iqlookfor{Replay limitation; a fix.}
\end{solution}

\begin{solution}{iq:mx:the-empirics-of-price-impact:5}
Part of the peak reverts after completion and part remains; how much depends on how informed the flow is, how fast liquidity returns and what else trades.
In the simulation all of it reverted, because nothing in it learns from the order; in real data a permanent part is usual.
\iqlookfor{Reversion; information; horizon.}
\end{solution}

\begin{solution}{iq:mx:the-empirics-of-price-impact:6}
Not necessarily: their orders are timed by forecasts, and the forecast drift is counted as impact. Compare after removing the forecast, or on orders sent
without one; the simulation's timed orders showed a 62\% higher prefactor with the same execution.
\iqlookfor{\omterm{def:mx:the-empirics-of-price-impact:alpha}{Alpha contamination}; a fair comparison.}
\end{solution}
